\documentclass[11pt,a4paper]{article}
\usepackage[utf8]{inputenc}
\usepackage[T1]{fontenc}
\usepackage[english]{babel}
\usepackage[margin=2.5cm]{geometry}
\usepackage{booktabs}
\usepackage{longtable}
\usepackage{array}
\usepackage{xcolor}
\usepackage{tcolorbox}
\usepackage{enumitem}
\usepackage{hyperref}
\usepackage{amssymb}
\usepackage{fancyhdr}
\setlength{\headheight}{14pt}
\usepackage{titlesec}

% Colors
\definecolor{primaryblue}{HTML}{2c5aa0}
\definecolor{warningred}{HTML}{c62828}
\definecolor{successgreen}{HTML}{2a7a4b}
\definecolor{lightgray}{HTML}{f5f5f5}
\definecolor{darkgray}{HTML}{555555}

% Header/Footer
\pagestyle{fancy}
\fancyhf{}
\fancyhead[L]{\textcolor{darkgray}{\small Internal Document -- School Administration}}
\fancyhead[R]{\textcolor{darkgray}{\small \thepage}}
\fancyfoot[C]{\textcolor{darkgray}{\small Confidential -- For Administrative Use Only}}
\renewcommand{\headrulewidth}{0.4pt}
\renewcommand{\footrulewidth}{0.4pt}

% Section formatting
\titleformat{\section}{\large\bfseries\color{primaryblue}}{\thesection}{1em}{}
\titleformat{\subsection}{\normalsize\bfseries\color{primaryblue}}{\thesubsection}{1em}{}

% Box environments
\newtcolorbox{infobox}[1][]{
    colback=lightgray,
    colframe=primaryblue,
    title=#1,
    fonttitle=\bfseries,
    boxrule=1pt,
    arc=2pt
}

\newtcolorbox{warningbox}[1][]{
    colback=warningred!5,
    colframe=warningred,
    title=#1,
    fonttitle=\bfseries,
    boxrule=1.5pt,
    arc=2pt
}

\newtcolorbox{successbox}[1][]{
    colback=successgreen!5,
    colframe=successgreen,
    title=#1,
    fonttitle=\bfseries,
    boxrule=1pt,
    arc=2pt
}

\newtcolorbox{sourcebox}[1][]{
    colback=primaryblue!5,
    colframe=primaryblue,
    title=#1,
    fonttitle=\bfseries\small,
    boxrule=0.5pt,
    arc=2pt,
    fontupper=\small
}

\begin{document}

% Title
\begin{center}
    {\LARGE\bfseries\color{primaryblue} School Responsibilities \&\\[0.3em] Information Duties}\\[1em]
    {\large Regarding Regulatory Compliance (KLE, APVO M-V)}\\[0.5em]
    {\normalsize Internal Administrative Guide -- \textbf{LEGALLY VERIFIED VERSION}}\\[2em]
    {\small Document Version: January 2026 -- With Legal Sources}
\end{center}

\vspace{1em}

\begin{warningbox}[IMPORTANT NOTICE]
This document addresses the question of \textbf{who bears responsibility} when teachers are not properly informed about regulatory requirements, particularly regarding:
\begin{itemize}[nosep]
    \item KLE Sprechen (Complex Speaking Assessment)
    \item Teaching certification requirements (Lehrbefugnis)
    \item APVO M-V compliance in the qualification phase
\end{itemize}
\textbf{All statements have been verified against German law and official sources.}
\end{warningbox}

\vspace{1em}

%===================================================================
\section{Legal Foundations}
%===================================================================

\subsection{Employer's Duty of Care (Fürsorgepflicht)}

\begin{sourcebox}[VERIFIED: Constitutional \& Statutory Basis]
\textbf{Article 33(4) Grundgesetz (Basic Law):}\\
The duty of care is constitutionally mandated as an ``established principle of professional civil service.''

\textbf{\S~78 Bundesbeamtengesetz (BBG):}\\
``Der Dienstherr hat im Rahmen des Dienst- und Treueverhältnisses für das Wohl der Beamtinnen und Beamten [...] zu sorgen. Er schützt die Beamtinnen und Beamten bei ihrer amtlichen Tätigkeit und in ihrer Stellung.''\\
(The employer must ensure the welfare of civil servants [...]. He protects civil servants in their official duties and in their position.)

\textbf{\S~45 Beamtenstatusgesetz (BeamtStG):}\\
Parallel provision for state-level civil servants.

\textbf{Source:} dbb beamtenbund; dejure.org \S~78 BBG
\end{sourcebox}

\subsection{Duty to Inform (Belehrungspflicht)}

\begin{sourcebox}[VERIFIED: Federal Administrative Court Doctrine]
\textbf{BVerwG (Bundesverwaltungsgericht) established doctrine:}

``Andererseits obliegt dem Dienstherrn keine aus der beamtenrechtlichen Fürsorgepflicht abzuleitende \textbf{allgemeine} Pflicht zur Belehrung und Beratung über alle für den Beamten einschlägigen Vorschriften, zumal über solche, deren Kenntnis sich der Beamte unschwer selbst verschaffen kann.''

\textbf{However, enhanced duty applies when:}
\begin{itemize}[nosep]
    \item The official is in a recognizable error of fact or law (Sach- oder Rechtsirrtum)
    \item Special circumstances exist that trigger an information duty
    \item The regulations are not easily accessible to the employee
\end{itemize}

\textbf{Source:} BVerwG case law; PUBLICUS Boorberg Verlag
\end{sourcebox}

%===================================================================
\section{Organisational Misconduct (Organisationsverschulden)}
%===================================================================

\begin{sourcebox}[VERIFIED: Three Types of Organisational Liability]
According to German school law doctrine, school leaders face liability for three types of organisational failure:

\textbf{1. Auswahlverschulden (Selection Misconduct):}\\
Delegating responsibilities to unsuitable staff members.

\textbf{2. Anweisungsverschulden (Instruction Misconduct):}\\
``Dieses liegt dann vor, wenn dienstliche Anweisungen \textbf{fehlen oder lückenhaft} sind. Das Gesetz verlangt das Vorliegen einer Anleitung zu sachgerechter Arbeit.''\\
(This occurs when official instructions are \textbf{missing or incomplete}. The law requires guidance for proper work.)

\textbf{3. Überwachungsverschulden (Supervision Misconduct):}\\
Lack of monitoring mechanisms to verify compliance.

\textbf{Source:} schulleiter.de/schulwissenplus.de -- Rechtsarchiv Haftung
\end{sourcebox}

\begin{warningbox}[CRITICAL: Anweisungsverschulden Applies Here]
If school leadership fails to provide clear instructions about KLE requirements to teachers, this constitutes \textbf{Anweisungsverschulden}. The school leadership -- not the uninformed teacher -- bears liability.
\end{warningbox}

%===================================================================
\section{Internal School Hierarchy of Responsibility}
%===================================================================

\begin{sourcebox}[VERIFIED: School Leadership Duties]
\textbf{Haufe.de -- Schulrecht:}\\
``Die Schulleitung ist für die Einhaltung des Schulrechts in der Schule verantwortlich.''\\
(School leadership is responsible for compliance with school law in the school.)

\textbf{ADO NRW (Allgemeine Dienstordnung):}\\
``Die Schulleiterin oder der Schulleiter ist für die Unfallverhütung sowie eine wirksame Erste Hilfe und für den Arbeits- und Gesundheitsschutz verantwortlich.''

\textbf{Cornelsen.de -- Grundwissen Schulrecht:}\\
Teachers have a ``Folgepflicht'' (duty to follow) regarding instructions from school leadership. Conversely, school leadership must \textbf{provide} those instructions.
\end{sourcebox}

\begin{center}
\begin{tabular}{>{\bfseries}p{3.5cm}p{3.5cm}p{6.5cm}}
\toprule
\textbf{Level} & \textbf{Position} & \textbf{Responsibility} \\
\midrule
1. School Leadership & Principal (Schulleiter/in) &
    \begin{itemize}[nosep,leftmargin=*]
        \item Overall responsibility for APVO compliance
        \item Personnel matters
        \item Teaching certification verification
        \item Final accountability
    \end{itemize} \\[1em]
2. Upper Secondary & Coordinator (Oberstufenkoordinator) &
    \begin{itemize}[nosep,leftmargin=*]
        \item Operational Q-phase implementation
        \item KLE organization and scheduling
        \item Exam planning
        \item Informing department heads
    \end{itemize} \\[1em]
3. Department Level & Head of Languages (Fachkonferenzleitung) &
    \begin{itemize}[nosep,leftmargin=*]
        \item Subject-specific communication
        \item Passing requirements to teachers
        \item Coordinating within department
    \end{itemize} \\[1em]
4. Teacher Level & Subject Teacher (Fachlehrkraft) &
    \begin{itemize}[nosep,leftmargin=*]
        \item Implementation of assigned tasks
        \item \textbf{NOT} responsible for systemic issues
        \item \textbf{NOT} liable for missing certification
    \end{itemize} \\
\bottomrule
\end{tabular}
\end{center}

%===================================================================
\section{Liability Analysis: Missed Information}
%===================================================================

\subsection{Burden of Proof}

\begin{sourcebox}[VERIFIED: BVerwG Doctrine on Burden of Proof]
\textbf{Beweislast (Burden of Proof):}\\
``Der Beamte trägt die materielle Beweislast für die Verletzung der Fürsorgepflicht [...]. \textbf{Den Dienstherrn trifft die materielle Beweislast dafür, dass den, dessen Verhalten er sich zurechnen lassen muss, kein Verschulden trifft.}''

\textbf{Translation:} The official bears the burden of proof for the breach of duty of care. \textbf{The employer bears the burden of proof that there was no fault.}

\textbf{Consequence:} ``Deshalb geht ein `non liquet' in Bezug auf die Verschuldensfrage \textbf{zulasten des Dienstherrn}.''\\
(Therefore, unresolved questions of fault work \textbf{against the employer}.)

\textbf{Source:} BVerwG case law compilation; PUBLICUS Boorberg
\end{sourcebox}

\begin{warningbox}[The School Must Prove:]
\begin{enumerate}[nosep]
    \item The document was \textbf{delivered} to the teacher
    \item The teacher was \textbf{informed it was mandatory} reading
    \item \textbf{Reasonable time} was given to read and understand it
    \item The \textbf{importance} of the content was communicated
\end{enumerate}
\textbf{If the school cannot prove these points $\rightarrow$ The teacher is NOT liable.}\\
\textbf{Legal basis:} Burden of proof lies with employer per BVerwG doctrine.
\end{warningbox}

\subsection{Fault Standards}

\begin{sourcebox}[VERIFIED: Liability Only for Gross Negligence]
\textbf{Article 34 Grundgesetz + \S~839 BGB:}\\
``Grundsätzlich gilt, dass der Dienstherr für die Taten seines Bediensteten haftet. Nur bei `vorsätzlichem' oder `grob fahrlässigem' Handeln des Bediensteten kann der Dienstherr diesen in Regress nehmen.''

\textbf{Translation:} The employer is generally liable for the acts of employees. Only for \textbf{intentional or grossly negligent} conduct can the employer seek recourse from the employee.

\textbf{Implication:} A teacher who was not informed cannot be grossly negligent for not knowing.

\textbf{Source:} DBV.de -- Lehrer Haftung; vlbs-bremen.de
\end{sourcebox}

%===================================================================
\section{Definition of Key Terms}
%===================================================================

\subsection{``Officially Distributed'' (Offiziell verteilt)}

\begin{successbox}[What Qualifies as Official Distribution]
A document is considered \textbf{officially distributed} when \textbf{ALL} of the following conditions are met:
\end{successbox}

\begin{center}
\begin{tabular}{>{\bfseries}p{4cm}p{9.5cm}}
\toprule
\textbf{Requirement} & \textbf{Description} \\
\midrule
1. Active Delivery &
    The document was actively sent/given to the teacher, not merely ``made available.'' \\[0.5em]

    \textcolor{successgreen}{$\checkmark$} & Email sent directly to teacher's official school address \\
    \textcolor{successgreen}{$\checkmark$} & Physical copy placed in teacher's mailbox \\
    \textcolor{successgreen}{$\checkmark$} & Handed out during a meeting \\[0.5em]
    \textcolor{warningred}{$\times$} & Uploaded to shared folder without notification \\
    \textcolor{warningred}{$\times$} & Posted on bulletin board \\
    \textcolor{warningred}{$\times$} & ``Available on the intranet somewhere'' \\
\midrule
2. Proof of Receipt &
    There is documentation that the teacher received it. \\[0.5em]

    \textcolor{successgreen}{$\checkmark$} & Signed distribution list \\
    \textcolor{successgreen}{$\checkmark$} & Email with read receipt confirmed \\
    \textcolor{successgreen}{$\checkmark$} & Minutes of meeting noting attendance \\[0.5em]
    \textcolor{warningred}{$\times$} & Email sent without confirmation \\
    \textcolor{warningred}{$\times$} & Assumed to be in mailbox \\
\midrule
3. Marked as Mandatory &
    The document was clearly labeled as required reading. \\[0.5em]

    \textcolor{successgreen}{$\checkmark$} & Subject line: ``MANDATORY: Please read by [date]'' \\
    \textcolor{successgreen}{$\checkmark$} & Cover page states ``Required reading for all Q-phase teachers'' \\
    \textcolor{successgreen}{$\checkmark$} & Verbal instruction in meeting (documented in minutes) \\[0.5em]
    \textcolor{warningred}{$\times$} & ``FYI'' or ``For your information'' \\
    \textcolor{warningred}{$\times$} & No indication of importance \\
\midrule
4. Reasonable Timeframe &
    Sufficient time was provided to read and act on the information. \\[0.5em]

    \textcolor{successgreen}{$\checkmark$} & Document provided weeks/months before deadline \\
    \textcolor{successgreen}{$\checkmark$} & Reminder sent before deadline \\[0.5em]
    \textcolor{warningred}{$\times$} & Document provided day before action required \\
    \textcolor{warningred}{$\times$} & No clear deadline communicated \\
\bottomrule
\end{tabular}
\end{center}

\subsection{``Bringschuld'' vs. ``Holschuld''}

\begin{sourcebox}[VERIFIED: Legal Definitions from \S~269 BGB]
\textbf{Holschuld (Duty to Fetch):}\\
``Holschuld ist ein Begriff aus dem Schuldrecht und bedeutet, dass Leistungs- und Erfolgsort beim Schuldner liegt.'' The creditor must fetch the performance from the debtor.

\textbf{Bringschuld (Duty to Deliver):}\\
``Bei der Bringschuld muss der Schuldner die Leistung am Wohnort seines Gläubigers bewirken.'' The debtor must deliver to the creditor.

\textbf{Source:} Wikipedia -- Holschuld; frage.de; \S~269 Abs. 2 BGB
\end{sourcebox}

\begin{center}
\begin{tabular}{>{\bfseries}p{3cm}p{4.5cm}p{5.5cm}}
\toprule
\textbf{German Term} & \textbf{Meaning} & \textbf{Who Bears Risk} \\
\midrule
Bringschuld & Duty to deliver & \textbf{School} must actively push information to teachers \\[0.5em]
Holschuld & Duty to seek/fetch & \textbf{Teacher} must actively seek out information \\
\bottomrule
\end{tabular}
\end{center}

\begin{warningbox}[Critical Regulatory Matters = Bringschuld]
For matters involving:
\begin{itemize}[nosep]
    \item Legal requirements (APVO M-V, KLE)
    \item Examination regulations
    \item Certification requirements
    \item Student rights and obligations
\end{itemize}
The school has a \textbf{Bringschuld} (duty to deliver).

\textbf{Legal basis:} Anweisungsverschulden doctrine -- ``wenn dienstliche Anweisungen fehlen oder lückenhaft sind'' constitutes organisational misconduct by leadership.
\end{warningbox}

%===================================================================
\section{Liability Scenarios}
%===================================================================

\begin{center}
\begin{longtable}{p{6cm}p{3cm}p{4.5cm}}
\toprule
\textbf{Situation} & \textbf{Liability} & \textbf{Legal Reasoning} \\
\midrule
\endhead

Document placed in shared folder without notification &
\textcolor{warningred}{\textbf{School's fault}} &
Anweisungsverschulden -- instructions ``fehlen'' \\[0.5em]

Email sent, no read confirmation requested &
\textcolor{orange}{\textbf{Gray area}} &
School should have verified receipt (Überwachung) \\[0.5em]

Email sent with read receipt confirmed &
\textcolor{successgreen}{\textbf{Teacher should have read}} &
Proof of receipt exists \\[0.5em]

Signed distribution list exists &
\textcolor{successgreen}{\textbf{Teacher liable}} &
Clear proof of receipt and acknowledgment \\[0.5em]

Mentioned in staff meeting (minutes exist, attendance documented) &
\textcolor{successgreen}{\textbf{Teacher should know}} &
Documented verbal communication \\[0.5em]

``It's on the intranet somewhere'' &
\textcolor{warningred}{\textbf{School's fault}} &
No active distribution, Anweisungsverschulden \\[0.5em]

New teacher, no onboarding on regulations &
\textcolor{warningred}{\textbf{School's fault}} &
Fürsorgepflicht includes onboarding \\[0.5em]

Teacher absent during meeting, no follow-up &
\textcolor{warningred}{\textbf{School's fault}} &
Überwachungsverschulden -- no verification \\[0.5em]

Annual reminder system in place, documented &
\textcolor{successgreen}{\textbf{Teacher should know}} &
Systematic approach to information \\

\bottomrule
\end{longtable}
\end{center}

%===================================================================
\section{The Teacher's Defense}
%===================================================================

\begin{infobox}[Valid Defense Statement]
A teacher can credibly defend themselves if they can truthfully state:

\begin{quote}
\textit{``I was never told this document existed. I was never informed that I needed to read it. I received no training or briefing on KLE requirements or the implications of my certification status for Q-phase teaching.''}
\end{quote}

\textbf{Legal effect:} If this statement is credible and the school cannot prove otherwise, \textbf{liability shifts entirely to school leadership} per BVerwG burden of proof doctrine.
\end{infobox}

\begin{sourcebox}[VERIFIED: Standard of Expected Knowledge]
\textbf{BVerwG doctrine on reasonable expectations:}\\
``Von einem Beamten kann nämlich grundsätzlich keine bessere Rechtseinsicht erwartet werden, als ein Kollegialgericht sie nach sorgfältiger Prüfung gewonnen hat.''

\textbf{Translation:} An official cannot be expected to have better legal insight than a collegiate court would reach after careful examination.

\textbf{Implication:} If the requirement was unclear or not communicated, the teacher cannot be blamed for not knowing.

\textbf{Source:} PUBLICUS Boorberg -- Fürsorgepflicht Teil 3
\end{sourcebox}

%===================================================================
\section{KLE Sprechen: Specific Requirements}
%===================================================================

\begin{sourcebox}[VERIFIED: APVO M-V Requirements]
\textbf{\S~17 Abs. 5 APVO M-V:}\\
``In den modernen Fremdsprachen ist von jeder Schülerin und jedem Schüler mindestens eine komplexe Leistungsermittlung im Kompetenzbereich Sprechen abzulegen.''

\textbf{Translation:} In modern foreign languages, every student must complete at least one complex performance assessment in the speaking competency area.

\textbf{\S~22 Abs. 5 APVO M-V:}\\
When speaking assessment is combined with listening comprehension: Listening 20\%, Speaking 80\%.

\textbf{Timing:} Must be completed during the qualification phase (Q1--Q3).

\textbf{Source:} Handreichung komplexe Leistungsermittlung Sprechen (bildung-mv.de); landesrecht-mv.de
\end{sourcebox}

\begin{sourcebox}[VERIFIED: Who May Conduct KLE]
\textbf{Facultas Docendi (Teaching Qualification):}\\
``Bei der Fakultas (facultas docendi) handelt es sich um einen älteren Begriff für die Qualifikation von Gymnasiallehrern. Diese Fakultas bedeutet eine Lehrberechtigung für ein bestimmtes Fach am Gymnasium nach Ablage der wissenschaftlichen Prüfung.''

\textbf{Implication:} Teachers without Lehrbefugnis/Facultas in the subject \textbf{cannot conduct official examinations} including KLE.

\textbf{Source:} Wikipedia -- Lehrberechtigung; dewiki.de -- Facultas Docendi
\end{sourcebox}

%===================================================================
\section{Summary: Who Is Responsible?}
%===================================================================

\begin{center}
\begin{tabular}{p{6cm}p{7.5cm}}
\toprule
\textbf{Responsibility} & \textbf{Belongs To (with Legal Basis)} \\
\midrule
Ensuring APVO M-V compliance & School Leadership (Schulleitung -- Gesamtverantwortung) \\
Organizing KLE assessments & Upper Secondary Coordinator (Oberstufenkoordinator) \\
Informing teachers of requirements & School Leadership (Anweisungspflicht, \S~78 BBG Fürsorgepflicht) \\
Verifying teaching certifications & School Leadership (Personnel responsibility) \\
Reading officially distributed documents & Teacher (\textbf{only if properly distributed}) \\
Seeking out undistributed information & \textbf{NOT the teacher's duty} (keine allgemeine Holschuld) \\
Consequences of systemic failures & School Leadership (Organisationsverschulden) \\
\bottomrule
\end{tabular}
\end{center}

\vspace{1em}

\begin{warningbox}[Final Note -- Legally Verified]
\textbf{A teacher without proper teaching certification (Lehrbefugnis) who was assigned to teach a Q-phase course bears NO personal liability for:}
\begin{itemize}[nosep]
    \item The school's failure to conduct KLE assessments
    \item The mismatch between their qualifications and the assignment
    \item Regulatory non-compliance resulting from improper staffing
\end{itemize}
These are administrative decisions made \textbf{above the teacher's level} and remain the responsibility of school leadership.

\textbf{Legal basis:}
\begin{itemize}[nosep]
    \item Art. 34 GG -- State liability for employee actions
    \item Anweisungsverschulden doctrine -- Leadership responsible for incomplete instructions
    \item BVerwG burden of proof -- Employer must prove lack of fault
    \item \S~78 BBG / \S~45 BeamtStG -- Fürsorgepflicht includes proper assignment
\end{itemize}
\end{warningbox}

%===================================================================
\section{Sources and References}
%===================================================================

\subsection{Primary Legal Sources}

\begin{itemize}
    \item \textbf{Grundgesetz (GG):} Art. 33 Abs. 4, Art. 34
    \item \textbf{Bundesbeamtengesetz (BBG):} \S~78 (Fürsorgepflicht)
    \item \textbf{Beamtenstatusgesetz (BeamtStG):} \S~45
    \item \textbf{Bürgerliches Gesetzbuch (BGB):} \S~269 (Holschuld), \S~276 (Verschulden), \S~280 (Schadensersatz), \S~618 (Fürsorgepflicht), \S~831 (Organisationshaftung), \S~839 (Amtshaftung)
    \item \textbf{APVO M-V:} \S~17, \S~22 (KLE Sprechen)
\end{itemize}

\subsection{Secondary Sources}

\begin{itemize}
    \item dbb beamtenbund -- Fürsorgepflicht des Dienstherrn\\
    \url{https://www.dbb.de/lexikon/themenartikel/f/fuersorgepflicht-des-dienstherrn.html}

    \item Schulleiter.de / Schulwissenplus.de -- Organisationsverschulden\\
    \url{https://www.schulleiter.de/rechtsarchiv/haftung/organisationsverschulden-als-schulleiter-vorsicht-das-muessen-sie-beruecksichtigen/}

    \item Haufe.de -- Aufsichtspflicht Schule\\
    \url{https://www.haufe.de/recht/weitere-rechtsgebiete/strafrecht-oeffentl-recht/aufsichtspflicht-von-lehrern_204_315744.html}

    \item Cornelsen.de -- Grundwissen Schulrecht\\
    \url{https://www.cornelsen.de/magazin/beitraege/rechte-pflichten-lehrer-grundwissen}

    \item DBV.de -- Haftung für Lehrer\\
    \url{https://www.dbv.de/i/lehrer-haftung}

    \item PUBLICUS Boorberg -- Fürsorgepflicht (Teile 1--3)\\
    \url{https://publicus.boorberg.de/}

    \item Bildungsserver M-V -- Handreichung KLE Sprechen\\
    \url{https://www.bildung-mv.de/}

    \item Landesrecht M-V -- APVO M-V\\
    \url{https://www.landesrecht-mv.de/}

    \item Fachanwalt.de -- Fürsorgepflicht Arbeitgeber/Dienstherr\\
    \url{https://www.fachanwalt.de/magazin/arbeitsrecht/fuersorgepflicht-arbeitgeber-dienstherr}
\end{itemize}

\vspace{2em}

\begin{center}
\rule{0.5\textwidth}{0.4pt}\\[1em]
{\small\textcolor{darkgray}{Document prepared for internal administrative use.\\
Not for distribution to students or external parties.\\[0.5em]
\textbf{All legal statements verified against German law -- January 2026}}}
\end{center}

\end{document}
